\subsection{Riesz 空间上的正线性形式}

我们回顾下述定义 (TVS, II, §2, No. 5)：

定义 1. --- 给定一个序向量空间 E，若 E 上的线性形式 L 满足 \(L(x) \geqslant 0\) 对 E 中每个 \(x \geqslant 0\) 成立，则称 L 为正的。

由于 \(L(y) - L(x) = L(y - x)\) ，这等价于说关系 \(x \leqslant y\) 蕴涵 \(L(x) \leqslant L(y)\) ，或者再说一遍，\(L\) 是 \(\mathbf{E}\) 上的递增函数。

例. --- 1) 设 A 是任意集，E 是定义在 A 上的一切实值函数所成的空间 \(R^{A}\) 的一个线性子空间。对每个元素 \(a \in A\) ，映射 \(x \mapsto x(a)\) 是 E 上的一个正线性形式。

\begin{enumerate}
\def\labelenumi{\arabic{enumi})}
\setcounter{enumi}{1}
\item
  设 \(\mathrm{I} = [a, b]\) 是 \(\mathbf{R}\) 中的一个紧区间，\(\mathrm{E}\) 是由 \(\mathrm{I}\) 上的正则实值函数组成的 Riesz 空间 (FRV, II, §1, No. 3)；映射 \(x \mapsto \int_{a}^{b} x(t) dt\) 是 \(\mathrm{E}\) 上的一个正线性形式。
\item
  设 F 是任意集，U 是 F 上的一个超滤子 (GT, I, §6, No. 4)，E 是 F 上有界实值函数所成的 Riesz 空间 \(\mathcal{B}(F)\) 。对每个 \(x \in E\) ，\(\lim_{\mathfrak{U}} x(t)\) 存在，因为 \(x(\mathfrak{U})\) 是相对紧集 \(x(F)\) 上一个超滤子的基，因而收敛。此外，若 \(x \geqslant 0\) ，则由不等式延拓原理，\(\lim_{\mathfrak{U}} x(t) \geqslant 0\) ；于是映射 \(x \mapsto \lim_{\mathfrak{U}} x\) 是 E 上的一个正线性形式。若把 U 取为由包含一个元素 \(a \in F\) 的集合组成的超滤子，就重新得到正线性形式 \(x \mapsto x(a)\) （例 1）。
\end{enumerate}

命题 1. --- 设 E 是序向量空间，L 是 E 到 R 中的一个映射，使得 \(L(x+y)=L(x)+L(y)\) ，且关系 \(x\geqslant0\) 蕴涵 \(L(x)\geqslant0\) ；则 \(L(\lambda x)=\lambda L(x)\) 对每个标量 \(\lambda\) 及每个 \(x\geqslant0\) 成立。

由于 \(L(-x) = -L(x)\) （L 是加群 E 在 R 中的一个表示），我们可以限于 \(\lambda \geqslant 0\) 的情形。对每个整数 \(n \geqslant 0\) ，我们有 \(L(nx) = nL(x)\) ，从而 \(L((1/n)x) = (1/n)L(x)\) ，因而 \(L(rx) = rL(x)\) 对每个有理数 \(r \geqslant 0\) 成立。另一方面，L 在 E 中递增；若 r 与 \(r'\) 是使 \(r \leqslant \lambda \leqslant r'\) 的有理数，则得 \(rL(x) \leqslant L(\lambda x) \leqslant r'L(x)\) ；由于 \(rL(x)\) 与 \(r'L(x)\) 可以任意接近 \(\lambda L(x)\) ，我们有 \(L(\lambda x) = \lambda L(x)\) 。

命题 2. --- 设 E 是实向量空间，C 是 E 中以 0 为顶点的凸锥且 E = C - C，又设 \(x \mapsto M(x)\) 是 C 到 R 中的一个映射，使得 \(M(\lambda x + \mu y) = \lambda M(x) + \mu M(y)\) 对一切 \(x \in C\) 、\(y \in C\) 、\(\lambda \geqslant 0\) 、\(\mu \geqslant 0\) 成立。则存在唯一的线性形式 L 把 M 延拓到 E 上。

由假设，每个 \(z \in E\) 都可写成 \(z = y - x\) ，其中 x, y 属于 C；此外，若又有 \(z = y' - x'\) ，其中 \(x' \in C\) ，\(y' \in C\) ，则

\(M(y) - M(x) = M(y') - M(x')\) ；因为由关系 \(y - x = y' - x'\) 推出 \(y + x' = x + y'\) ，从而 \(M(y) + M(x') = M(x) + M(y')\) 。我们把 \(L(z)\) 定义为 \(M(y) - M(x)\) 当 \(z\) 表为 C 的两元素之差 \(y - x\) 时所取的公共值；立即可以验证，\(L\) 是 E 上延拓 \(M\) 的线性形式；\(L\) 的唯一性由 C 生成空间 E 这一事实得出。

命题 3. --- 设 E 是有向序向量空间，P 是 E 中元素 \(\geqslant 0\) 所成的集，又设 \(x \mapsto M(x)\) 是 P 到 R 中的映射，取值 \(\geqslant 0\) ，使得对 P 中一切 x, y 有 \(M(x + y) = M(x) + M(y)\) 。则存在唯一的正线性形式 L 把 M 延拓到 E 上。

由于 E = P - P，与命题 2 中同样的论证首先证明：存在唯一的 E 到 R 中的加性映射 L 延拓 M。于是命题 1 表明，\(L(\lambda x) = \lambda L(x)\) 对一切 \(\lambda \geqslant 0\) 及一切 \(x \in P\) 成立，由此立即可知 L 是线性形式。

